Esta tabla es el mapa de toda la réplica. Las columnas son las formas de correr las pruebas; las filas, los experimentos.

En el Ultra 9 hay cuatro tandas grandes: los binarios oficiales con Docker y sin Docker, el programa del equipo sin Docker, y una copia de ese programa con una corrección de una palabra que explicaremos en la parte 6.

La penúltima columna es importante. El guion del equipo, ejecutar\_pruebas.sh, se corrió igual en las dos máquinas. Eso nos da una comparación directa entre el i5 y el Ultra 9 con el mismo programa y los mismos parámetros. La veremos en la parte 5.
